\scanpage{314}
\mspage{264}

en changeant les rôles de $X_1$ et $X_2$, on voit qu'on peut supposer de plus que $S\mathbin{\unsure{-|\to}}X_2$ est également imm. \inserted{ouv.} de Dwyer.

Considérons alors

\diagram{\(\begin{tikzcd}[ampersand replacement=\&,column sep=large,row sep=large]
S\arrow[r,arrows={Hooks[right]-}]\arrow[d,hook]\arrow[dr,phantom,"{Q'}" description]\&X_1\arrow[d,hook]\&{}\\
X_2\arrow[r,arrows={Hooks[right]-}]\&X_1\amalg_S X_2\arrow[r,"{\unsure{\pi|\Pi}}"]\&X_3\\
{}\&\unsure{\mathscr X|\mathcal X}_Q\arrow[u,"{\varphi_{Q'}}"]\arrow[ur,"{\varphi_Q}"']\&{}
\end{tikzcd}.\)}

On a $\varphi_Q=\pi\varphi_{Q'}$, d'autre part on sait que $\varphi_{Q'}\in W$ (cor. 3 p. 149). Donc

\[
\pi\in W\Longleftrightarrow\varphi_Q\in W,
\]

or $\pi\in W$ signifie que $Q$ est $W$-cocartésien, qed.
% archive sheet number: 314
